\chapter{Temporal-Difference Learning}

Chapters 3 and 4 already give two ways to compute values. Neither is ideal in every setting. Dynamic Programming needs the model $P$ and a full backup over every next pair $(s',r)$; Monte Carlo needs no model and uses episode returns $G_t$. Temporal-Difference (TD) learning (\citetxt{sutton1988}{Sutton, 1988}) sits between them: it keeps Monte Carlo's model-free samples and Dynamic Programming's bootstrapping, so the table can be updated from one observed pair $(R_{t+1},S_{t+1})$ rather than waiting for $G_t$. The algorithms below evaluate a given $\pi$ and return $V\approx V^\pi$. They do not replace $\pi$: that Improve step was Section 4.4 for Monte Carlo, and is Chapter 6 for TD.

Those DP limitations follow from the backup: it is a sum over $P$. In Dynamic Programming (Chapter 3), the model $P(s',r\mid s,a)$ is known, so each update is a full backup. Value iteration and the improve step of policy iteration use

\[
V(s)
\leftarrow
\max_a
\sum_{s',r}
P(s',r\mid s,a)\bigl[r+\gamma V(s')\bigr],
\]

while policy evaluation uses the same sum with the fixed action $a=\pi(s)$:

\[
V(s)
\leftarrow
\sum_{s',r}
P(s',r\mid s,\pi(s))\bigl[r+\gamma V(s')\bigr].
\]

Every backup averages over all next pairs $(s',r)$ under $P$, and replaces the rest of the return by the current table entry $V(s')$---that reuse of $V$ is \textbf{bootstrapping}. No episodes are required, only sweeps over the state table. The price is dependence on $P$: if the model is unknown or the sum is infeasible, DP cannot run.

When $P$ is unavailable, Monte Carlo (Chapter 4) generates episodes under $\pi$ and never uses $P$. The return obeys (1.10), $G_t=R_{t+1}+\gamma G_{t+1}$, so $V^\pi(s)=\mathbb{E}_\pi[G_t\mid S_t=s]=\mathbb{E}_\pi[R_{t+1}+\gamma G_{t+1}\mid S_t=s]$. The average of returns from visits to $s$ is

\[
V^\pi(s)
\approx
\frac{1}{N}\sum_{i=1}^{N} G^{(i)}(s),
\]

and the next return is folded in with step-size $\alpha$ as in (4.3)---the \textbf{Monte Carlo} update

\[
V(S_t)
\leftarrow
V(S_t)+\alpha\bigl(G_t-V(S_t)\bigr).
\]

(Here $\alpha=1/(N{+}1)$ recovers an exact sample mean after $N$ prior samples for that state.) Because $G_t$ is built only from observed rewards, there is no sum over $P$ and no lookup of $V(S_{t+1})$. Two difficulties remain. First, $G_t$ stacks many random rewards, so the return from a given state has high variance across visits. Second, $G_t$ is known only when the episode ends, so updates wait for termination: long episodes delay learning, and continuing tasks do not fit.

Both difficulties come from using the full return as the learning target. Equivalently, the same $V^\pi$ takes the compact form, which replaces $G_{t+1}$ by $V^\pi(S_{t+1})$ (\eqref{eq:B.11} in \appref{B}):

\begin{equation}
V^\pi(s)
=
\mathbb{E}_\pi\bigl[R_{t+1}+\gamma V^\pi(S_{t+1})\mid S_t=s\bigr].
\tag{5.1}
\end{equation}

Run the same sample-mean steps as for Monte Carlo---average the one-step backup over visits, then step with $\alpha$---but replace the unknown $V^\pi(S_{t+1})$ by the current table $V(S_{t+1})$. The result is the \textbf{TD} update

\begin{equation}
V(S_t) \leftarrow V(S_t) + \alpha \left[ R_{t+1} + \gamma V(S_{t+1}) - V(S_t) \right]
\tag{5.2}
\end{equation}

Like Monte Carlo, (5.2) needs no $P$: the pair $(R_{t+1},S_{t+1})$ is observed from interaction. Like Dynamic Programming, it bootstraps with a table lookup at the next state. Replacing $G_{t+1}$ by $V(S_{t+1})$ avoids waiting for the episode to end and avoids the high-variance full return $G_t$. The backup $R_{t+1}+\gamma V(S_{t+1})$ still varies across visits, because $R_{t+1}$ and $S_{t+1}$ are random and $V$ is not $V^\pi$. Section 5.1 develops this \textbf{one-step TD} method (TD error and algorithm); Sections 5.2--5.3 extend it to $n$-step TD and TD($\lambda$). Control---$Q$ and an improving $\pi$---is Chapter 6.

\section{One-step TD (TD(0))}

The update (5.2) steps $V(S_t)$ toward the backup $R_{t+1}+\gamma V(S_{t+1})$. That backup uses only the next reward and the table at $S_{t+1}$---a \textbf{one-step} backup. The difference between that backup and the prediction is the \textbf{TD error}---the signal that tells whether $V(S_t)$ should rise or fall:

\begin{equation}
\delta_t
=
\bigl(R_{t+1}+\gamma V(S_{t+1})\bigr) - V(S_t).
\tag{5.3}
\end{equation}

Equivalently, (5.2) is

\[
V(S_t)
\leftarrow
V(S_t)+\alpha\,\delta_t.
\]

If $\delta_t>0$, then $R_{t+1}+\gamma V(S_{t+1})>V(S_t)$, so $V(S_t)$ increases; if $\delta_t<0$, it decreases. That push and pull of $V$ by $\delta_t$ is the learning rule. The lookup $V(S_{t+1})$ in the backup is the bootstrap; $V(S_t)$ is only the prediction being corrected. Both come from the table, so $V$ must be initialized before the first step (an arbitrary table; a common choice is $V\equiv 0$). If $S_{t+1}$ is terminal, take $V(S_{t+1})=0$ and the backup is just $R_{t+1}$.

The step in (5.2) is not the Monte Carlo $1/n$ of \appref{D}. That derivation averages $G$; with $\pi$ and $P$ fixed, $\mathbb{E}[G]=V^\pi$ does not depend on the table. Here the target is $R_{t+1}+\gamma V(S_{t+1})$, which moves with $V$ even when $\pi$ does not. Early backups use a poor $V$; later backups use a better $V$. A $1/n$ weight would fall on the early backups. A constant $\alpha\in(0,1]$ keeps weight on later backups as $V$ improves, and is the usual practical choice.

Unlike Monte Carlo, TD does \textbf{not} generate a full episode and then update. The outer loop starts another episode; the inner loop generates that episode under $\pi$ and applies (5.2) to each observed pair $(R_{t+1},S_{t+1})$. That one-step method is also called \textbf{TD(0)} (the name is explained in Section 5.3):

\begin{implnote}[Algorithm (one-step TD prediction, TD(0))]
Input: policy $\pi$ to evaluate.

Output: estimate $V\approx V^\pi$.

\begin{enumerate}
\item Initialize $V(s)$ arbitrarily (e.g. $0$), with $V(\text{terminal})=0$.
\item Loop for each episode:
\begin{itemize}
\item Initialize $S_0$.
\item For $t=0,1,2,\ldots$ until termination:
\begin{itemize}
\item Take $A_t\sim\pi(\cdot\mid S_t)$, observe $R_{t+1}$, $S_{t+1}$.
\item $V(S_t)\leftarrow V(S_t)+\alpha\bigl[R_{t+1}+\gamma V(S_{t+1})-V(S_t)\bigr]$.
\item $S_t\leftarrow S_{t+1}$.
\end{itemize}
\end{itemize}
\item After enough episodes, return $V$ ($\approx V^\pi$).
\end{enumerate}

\end{implnote}

\lstinputlisting{../code/ch05/td_prediction.py}


\section{$n$-step TD}

TD(0) steps toward one reward and then the table at the next state. A reward that arrives several steps later can change $V(S_t)$ only by passing through those later table entries, and if $V$ is still wrong those entries pollute the backup. $n$-step TD waits for the next $n$ rewards and bootstraps only at $S_{t+n}$, so those rewards enter the target directly. The case $n=1$ is TD(0); waiting until termination recovers the Monte Carlo return $G$. Larger $n$ uses more observed rewards and less of $V$, so less bootstrap error and faster credit, but more variance and a delay of $n$ steps. Intermediate $n$ is often better than either extreme.

The one-step backup already in (5.2) is

\[
U_t^{(1)}=R_{t+1}+\gamma V(S_{t+1}).
\]

Two rewards and a bootstrap at $S_{t+2}$ start from the return identity (1.10), $G_t=R_{t+1}+\gamma G_{t+1}$. Substitute $G_{t+1}=R_{t+2}+\gamma G_{t+2}$:

\begin{equation}
G_t=R_{t+1}+\gamma R_{t+2}+\gamma^2 G_{t+2}.
\tag{5.5}
\end{equation}

$V^\pi(s)=\mathbb{E}_\pi[G_t\mid S_t=s]$. Take that expectation on (5.5) and replace $G_{t+2}$ by $V^\pi(S_{t+2})$ (tower property, as in (5.1)):

\[
V^\pi(s)
=
\mathbb{E}_\pi\bigl[R_{t+1}+\gamma R_{t+2}+\gamma^2 V^\pi(S_{t+2})\mid S_t=s\bigr].
\]

One sample $(R_{t+1},R_{t+2},S_{t+2})$ with the table $V$ in place of $V^\pi$ is the two-step backup (write $U$ for a backup; $G$ stays the true return)

\[
U_t^{(2)}=R_{t+1}+\gamma R_{t+2}+\gamma^2 V(S_{t+2}).
\]

When $V=V^\pi$, $\mathbb{E}_\pi[U_t^{(2)}\mid S_t=s]=V^\pi(s)$. The update is (5.2) with this target:

\[
V(S_t)\leftarrow V(S_t)+\alpha\bigl(U_t^{(2)}-V(S_t)\bigr).
\]

The same pattern with $n$ rewards and a bootstrap at $S_{t+n}$ is

\[
U_t^{(n)}=R_{t+1}+\gamma R_{t+2}+\cdots+\gamma^{n-1}R_{t+n}+\gamma^n V(S_{t+n}),
\]

\[
V(S_t)\leftarrow V(S_t)+\alpha\bigl(U_t^{(n)}-V(S_t)\bigr).
\]

For $n=1$, $U_t^{(1)}$ is the backup in (5.2). If the episode ends at or before $t+n$, then $V(\text{terminal})=0$ and $U_t^{(n)}$ is the remaining return $G_t$. $U_t^{(n)}$ contains $V$, so a constant $\alpha$ is the same practical choice as in Section 5.1.

The method that uses a fixed $n$ is \textbf{$n$-step TD}. It needs the next $n$ transitions (or termination), so an episode is generated first, then $V$ is updated with $U_t^{(n)}$ from a frozen table.

\begin{implnote}[Algorithm ($n$-step TD prediction)]
Input: policy $\pi$ to evaluate; horizon $n\ge 1$.

Output: estimate $V\approx V^\pi$.

\begin{enumerate}
\item Initialize $V(s)$ arbitrarily (e.g. $0$), with $V(\text{terminal})=0$.
\item Loop for each episode:
\begin{itemize}
\item Generate an episode following $\pi$:   $S_0,A_0,R_1,\;\ldots,\;S_{T-1},A_{T-1},R_T$.
\item For each $t=0,1,\ldots,T{-}1$, compute $U_t^{(n)}$ (same frozen table $V$ for every $t$).
\item For each $t=0,1,\ldots,T{-}1$:   $V(S_t)\leftarrow V(S_t)+\alpha\bigl(U_t^{(n)}-V(S_t)\bigr)$.
\end{itemize}
\item After enough episodes, return $V$ ($\approx V^\pi$).
\end{enumerate}

\end{implnote}

\lstinputlisting{../code/ch05/n_step_td_prediction.py}


Section 5.3 mixes several horizons $n$ into one backup and obtains \textbf{TD($\lambda$)}.

\section{TD($\lambda$)}

\textbf{Mixing backups.} In Section 5.2 the learner must pick one fixed horizon $n$: use only $U_t^{(n)}$. That commits to a single tradeoff (short $n$ bootstraps early; long $n$ waits for more rewards). A single horizon is not required. One can instead form \textbf{one} backup by averaging several horizons, then step $V$ toward that mix as in (4.3). For example, give equal weight to the 1-step, 2-step, and 3-step backups:

\[
U
=
\frac{1}{3}U_t^{(1)}+\frac{1}{3}U_t^{(2)}+\frac{1}{3}U_t^{(3)}
=
\sum_{n=1}^{3}\frac{1}{3}\,U_t^{(n)}.
\]

Then

\[
V(S_t)\leftarrow V(S_t)+\alpha\bigl(U-V(S_t)\bigr)
=
V(S_t)+\alpha\Biggl(\sum_{n=1}^{3}\frac{1}{3}\,U_t^{(n)}-V(S_t)\Biggr).
\]

More generally, pick any nonnegative weights $w_n$ that sum to one ($\sum_n w_n=1$) and set

\[
U=\sum_{n} w_n\,U_t^{(n)}.
\]

The equal mix above is the special case $w_1=w_2=w_3=1/3$. The update is again

\[
V(S_t)\leftarrow V(S_t)+\alpha\Biggl(\sum_{n} w_n\,U_t^{(n)}-V(S_t)\Biggr).
\]

The weights must sum to $1$ so $U$ stays a weighted \textbf{average} on the scale of a return (and of $V$). Uniform $w_n=1/N$ is fine for a \textbf{finite} list $n=1,\ldots,N$. It does \textbf{not} extend to $N\to\infty$: you cannot give every $n=1,2,3,\ldots$ the same positive weight and still have $\sum_{n=1}^{\infty} w_n=1$. A flat $1/N$ also gives the same importance to a short backup (few rewards, early use of $V$) and a long backup (many random rewards). Usually one wants short horizons to count more, and long horizons less---that leads to the geometric weights next.

\textbf{Geometric weights and $\lambda$.} Let each longer backup get a fixed fraction $\lambda\in[0,1)$ of the previous weight. The unnormalized weights are $1,\lambda,\lambda^2,\ldots$, i.e. $\lambda^{n-1}$. The first $N$ sum to

\[
1+\lambda+\cdots+\lambda^{N-1}=\frac{1-\lambda^N}{1-\lambda}\qquad(\lambda\neq 1).
\]

If $N\to\infty$ and $\lambda<1$, then $\lambda^N\to 0$, so the infinite sum is $1/(1-\lambda)$. Normalizing gives

\[
w_n=(1-\lambda)\,\lambda^{n-1}.
\]

As $n$ grows, $\lambda^{n-1}$ shrinks ($\lambda<1$), so $w_n$ shrinks. A short backup $U_t^{(1)}$ (one reward, then $V$) has the largest weight; a long backup $U_t^{(n)}$ (many rewards, later $V$) has less. For $\lambda=0.5$: $w_1=1/2$, $w_2=1/4$, $w_3=1/8$, $\ldots$. The resulting mix is the \textbf{$\lambda$-backup}

\begin{equation}
U_t^{\lambda}
=
\sum_{n=1}^{\infty} w_n\,U_t^{(n)}
=
(1-\lambda)\sum_{n=1}^{\infty}\lambda^{n-1}U_t^{(n)},
\tag{5.4}
\end{equation}

and

\[
V(S_t)\leftarrow V(S_t)+\alpha\bigl(U_t^{\lambda}-V(S_t)\bigr).
\]

That method is \textbf{TD($\lambda$)}. At $\lambda=0$, $U_t^{\lambda}=U_t^{(1)}$, which is the one-step backup of Section 5.1 and the $n{=}1$ case of Section 5.2, so that method is named \textbf{TD(0)}. As $\lambda\to 1$, $U_t^{\lambda}\to G_t$ (Monte Carlo)---\textbf{TD(1)}.

\begin{asidenote}[Note]
Equation (5.4) takes $\lambda\in[0,1)$. For every such $\lambda$ the weights still sum to $1$,
\[
\sum_{n=1}^{\infty}(1-\lambda)\lambda^{n-1}=1,
\]
so $U_t^{\lambda}$ is an average, not $0$. Substituting $\lambda=1$ into (5.4) is invalid ($1-\lambda=0$). TD(1) is $\lambda\to 1$ of that average.

Unrolling (1.10) $n$ times (the same pattern as (5.5)) gives
\[
G_t=R_{t+1}+\gamma R_{t+2}+\cdots+\gamma^{n-1}R_{t+n}+\gamma^n G_{t+n}.
\]
Section 5.2 already has
\[
U_t^{(n)}=R_{t+1}+\gamma R_{t+2}+\cdots+\gamma^{n-1}R_{t+n}+\gamma^n V(S_{t+n}).
\]
The first $n$ rewards match, so
\[
U_t^{(n)}-G_t=\gamma^n\bigl(V(S_{t+n})-G_{t+n}\bigr).
\]
$G$ is bounded as in Chapter 1: with $\gamma\in[0,1)$ and bounded rewards, (1.8) converges. $V$ is a finite table, frozen while this backup is formed. With $\alpha\in(0,1]$ the update $V\leftarrow(1-\alpha)V+\alpha U$ is a convex combination; for one-step TD, $U=R+\gamma V(s')$, so if the table already lies in the same interval that bounds $G$, then $U$ stays in that interval and so does $V$. The factor $\gamma^n$ therefore sends the difference to $0$, hence $U_t^{(n)}\to G_t$ as $n\to\infty$. As $\lambda\to 1$, each fixed-$n$ weight vanishes, so the average in (5.4) is carried by large $n$, where $U_t^{(n)}$ is already near $G_t$. Thus $U_t^{\lambda}\to G_t$. The update becomes
\[
V(S_t)\leftarrow V(S_t)+\alpha\bigl(G_t-V(S_t)\bigr),
\]
which is the Monte Carlo step (4.3).
\end{asidenote}

Here $\lambda$ is not a power on $\gamma$ and not the horizon $n$: $n$-step TD (Section 5.2) picks one $U_t^{(n)}$, while $\lambda$ sets how fast the mix forgets longer backups. There is no TD(2) in the $\lambda$ naming. From here on, Section 5.1's method is called TD(0) when contrasting with multi-step or $\lambda$ methods.

On a finite episode that ends at time $T$, the non-terminal states are $S_0,\ldots,S_{T-1}$. From a time $t$ in that episode, Section 5.2's backup $U_t^{(n)}$ uses the next $n$ rewards and then $V(S_{t+n})$. Here $S_{t+n}$ is the state $n$ steps after $S_t$ on the same episode. That state is a non-terminal only while $t+n<T$. If $n\ge T{-}t$, the episode has already ended, so Section 5.2 gives $U_t^{(n)}=G_t$ (the lookup $V(\text{terminal})$ is $0$). The mix (5.4) is unchanged: every term with $n\ge T{-}t$ equals $G_t$. The weights still sum to $1$. The first $T{-}t{-}1$ of them add to $1-\lambda^{T-t-1}$, and the leftover $\lambda^{T-t-1}$ therefore sits on $G_t$:
\begin{equation}
U_t^{\lambda}
=
\sum_{n=1}^{T-t-1}(1-\lambda)\lambda^{n-1}U_t^{(n)}
+
\lambda^{T-t-1}G_t.
\tag{5.6}
\end{equation}
Computing $U_t^{\lambda}$ needs the future rewards of the episode, so the forward-view algorithm below generates an episode first, then updates toward (5.6)---unlike online TD(0).

\begin{implnote}[Algorithm (forward-view TD($\lambda$) prediction)]
Input: policy $\pi$ to evaluate; $\lambda\in[0,1]$.

Output: estimate $V\approx V^\pi$.

\begin{enumerate}
\item Initialize $V(s)$ arbitrarily (e.g. $0$), with $V(\text{terminal})=0$.
\item Loop for each episode:
\begin{itemize}
\item Generate an episode following $\pi$:   $S_0,A_0,R_1,\;S_1,A_1,R_2,\;\ldots,\;S_{T-1},A_{T-1},R_T$.
\item For each $t=0,1,\ldots,T{-}1$, compute $U_t^{\lambda}$ from (5.6) (same frozen table $V$ for every $t$).
\item For each $t=0,1,\ldots,T{-}1$:   $V(S_t)\leftarrow V(S_t)+\alpha\bigl(U_t^{\lambda}-V(S_t)\bigr)$.
\end{itemize}
\item After enough episodes, return $V$ ($\approx V^\pi$).
\end{enumerate}

\end{implnote}

\lstinputlisting{../code/ch05/td_lambda_prediction.py}

